\subsubsection{An asset floor changes the allocation of a unit's loss}
\label{sec:impairment-floors}

A cash-generating unit can lose value while one of its assets retains a
substantial disposal value. Allocating the unit's loss in proportion to asset
balances just once may then reduce that asset too far. IAS 36.104--105 first
allocates loss to goodwill, then to the other assets pro rata, while protecting
each asset from a reduction below the highest of zero, its fair value less
costs of disposal if measurable, and its value in use if determinable. A share
that cannot be assigned to one asset is redistributed to the others
\citep{eu20231803}.

Consider a synthetic cost-model unit containing equipment A with carrying
amount 60, equipment B with carrying amount 40, and allocated acquired
goodwill 20. Cash 80 lies outside the tested unit. With no liabilities, opening
equity is 200. The unit boundary and goodwill allocation are assumed to have
been admitted under paragraphs 66--80; there are no shared corporate assets,
noncontrolling interests or liabilities requiring an adjustment to this
comparison. In a real case those judgments precede the calculation. In
particular, paragraph 78 requires consistent treatment of a liability that is
included in the unit's recoverable value; subtracting it from only the carrying
amount would compare different populations.

At measurement date $T_0$, admitted value in use is 70 and fair value less
disposal costs is 65, so the unit's recoverable amount is 70. Individual
values in use are supported as not determinable, while the individual disposal
values are 55 for A and 5 for B. These are supplied valuation assumptions,
not values inferred by an allocation routine. They establish floors of 55
and 5. A missing valuation would be a different state: the program may omit
a term because its non-determinability is admitted, but may not replace an
unexamined input with zero.

The unit's loss is $60+40+20-70=50$. Goodwill takes the first 20. An
unconstrained allocation of the remaining 30 would assign
$30\times60/100=18$ to A and $30\times40/100=12$ to B. A, however, can
absorb only $60-55=5$. Its excess share of 13 passes to B, whose total loss
becomes $12+13=25$, within its capacity $40-5=35$. The closing balances are
therefore
\[
  A=60-5=55,\qquad B=40-25=15,\qquad
  \text{goodwill}=20-20=0.
\]
Together they equal the unit's recoverable amount of 70. Debit impairment
expense 50, credit goodwill 20, and credit the separate asset impairment
allowances by 5 and 25. Cash 80 plus equipment 70 equals equity 150. The
unconstrained 18/12 entry would also balance, but would leave A at 42,
below its admitted floor of 55. Clipping A's loss to 5 without redistributing
the excess would instead recognize only 17 of the required remaining 30.

For an implementer, each round has a remaining amount and an active set of
assets. Compute shares from their pre-allocation carrying amounts. Fix each
share that reaches an asset's available capacity, remove that asset from the
active set, and redistribute the remaining amount. A round that needs further
redistribution removes at least one asset, so the procedure terminates after
finitely many such removals. For positive carrying weights $b_i$ and
capacities $c_i$, its allocation has the form
\[
 x_i=\min(c_i,\lambda b_i),\qquad
 \sum_i x_i=\min\!\left(q,\sum_i c_i\right),
\]
where $q$ is the requested allocation and $\lambda\geq0$ gives the common
proportional rate among assets that have not reached their limits. In the
example A is fixed at 5 and B receives the remaining 25. A source floor above
an asset's existing cost gives zero reduction capacity; it does not create
permission to increase that asset during an impairment calculation.

The capacity limit also explains how to handle an unsuccessful allocation.
A diagnostic input requesting 50 against capacities 5 and 35 can allocate
only 40, leaving 10. That arithmetic does not establish that the supplied
unit and individual valuations are jointly coherent. It calls for checking
measurement assumptions, the unit boundary and any separate obligations.
Paragraph 108 permits a liability for a remaining unit loss only if another
IFRS requires one. AccDSL must retain the residual and the unresolved decision;
it cannot invent a payable to make the unit reach a requested total. This
diagnostic is not a second admitted company example.

\subsubsection{A recovery is limited by the depreciation that would have occurred}
\label{sec:impairment-reversal-memory}

Suppose both pieces of equipment have five remaining years of even use and
zero residual values at $T_0$. Under paragraph 63, their revised carrying
amounts give next-year depreciation $55/5=11$ and $15/5=3$. At $T_1$, one
full year later, they carry 44 and 12. The company must also retain the path
that would have occurred without impairment. On that path depreciation would
have been $60/5=12$ and $40/5=8$, leaving 48 and 32. Restoring a past loss
therefore cannot mean restoring the old nominal balances.

Assume a supported change in expected future cash flows now gives the unit a
recoverable amount of 90. Individual recoverable amounts remain admitted as
not determinable. Paragraphs 114--116 require the change in estimates; mere
passage of time and discount unwinding would not support this reversal.
Paragraphs 122--123 allocate a unit reversal over assets other than goodwill
pro rata, subject to each asset's recoverable amount if determinable and its
depreciated no-impairment amount. Goodwill itself cannot be restored under
paragraph 124 \citep{eu20231803}.

The requested increase is $90-(44+12)=34$. A has room for only $48-44=4$;
B has room for $32-12=20$. The same redistribution procedure therefore
recognises $4+20=24$, leaving equipment at 48 and 32 and goodwill at zero.
The recognized total of 80 is below the unit's value of 90. That difference
does not identify an omitted asset, authorize goodwill recognition or imply
cash has been received.

Debit the impairment allowances by 4 and 20 and credit reversal gain 24.
The original allowances of 5 and 25 consequently retain balances of 1 and 5.
Accumulated depreciation is 11 and 3. The register now reconciles as
\[
  \text{PPE}=100-\underbrace{(11+3)}_{\text{depreciation}}
                 -\underbrace{(1+5)}_{\text{impairment allowances}}=80.
\]
The remaining allowances record the different past depreciation paths.
Erasing all 30 of the original allowances would produce PPE of 86, exceeding
the no-impairment ceiling by 6. Future annual depreciation is $48/4=12$ and
$32/4=8$, consistent with the remaining lives.

\begin{center}
\begin{tabular}{@{}lrr@{}}
\toprule
Pretax amount & $T_0$ impairment & Following year to $T_1$\\
\midrule
Impairment expense & 50 & 0\\
Depreciation / reversal gain & 0 / 0 & 14 / 24\\
Profit / OCI & $-50$ / 0 & 10 / 0\\
Closing cash / PPE & 80 / 70 & 80 / 80\\
Closing goodwill / liabilities & 0 / 0 & 0 / 0\\
Closing equity & 150 & 160\\
Operating cash & 0 & 0\\
\bottomrule
\end{tabular}
\end{center}

The first column isolates the impairment measurement; the second follows the
next full year. With no cash or owner movements in this slice, the operating
cash reconciliations are $-50+50=0$ and $10+14-24=0$, removing the noncash
loss, depreciation and reversal under IAS 7.20. Equity moves from 200 to
150, then to 160. Tax remains a separate calculation under IAS 36.64 and
IAS 12. These technical rules use the retained original 2023 EU text, with
its date corrigendum; the relative dates do not assert current reporting
admission or a complete later amendment review.

The method thus needs actual and no-impairment depreciation histories,
measurement dates, individual valuation-status judgments, goodwill identity
and evidence of the changed estimate. There is also an explicit numerical
boundary. If every eligible asset has zero carrying amount but a positive
reversal is requested, the positive-weight procedure above has no denominator.
It returns an unresolved allocation basis and preserves the history; it does
not choose equal weights without an additional justified interpretation.
That extension, along with shared corporate assets and other scope
interactions, remains distinct from the fully specified positive-weight case.
